\section{\texorpdfstring{Admissible subcategories:\allowbreak{} hypotheses,\allowbreak{} projections,\allowbreak{} and receiving criteria}{Admissible    and receiving criteria}}\label{proofs-11665-12003-admissible-subcategories-hypotheses-projections-and-receiving-criteria}

Authority:\allowbreak{} derived.tex at a04446e5\allowbreak{}7ec1fbc2\allowbreak{}52a871af\allowbreak{}cec7752f\allowbreak{}b2807b14, original lines 11665–\allowbreak{}12003. Current-source comparisons use the frozen public preimages at 95a51593\allowbreak{}aceb247a\allowbreak{}c7a18311\allowbreak{}1678c3a0\allowbreak{}f8b8f4b5. This note supports PART21. The source cites Bondal–\allowbreak{}Kapranov,\allowbreak{} §1; the arguments below are proved from the specified triangulated-category axioms and the earlier source lemmas,\allowbreak{} without claiming to have read that reference in this pass or claiming novelty.

\subsection{1. Orthogonals and the two exact Hom criteria}\label{proofs-11665-12003-orthogonals-and-the-two-exact-hom-criteria}

For a full subcategory A of an additive category D,\allowbreak{} the original right orthogonal A\^{}⊥ has objects Y for which Hom\_\allowbreak{}D(A,\allowbreak{}Y)\allowbreak{}\allowbreak{}=0 for every A\allowbreak{}∈A. The original left orthogonal \^{}⊥A has objects Y for which Hom\_\allowbreak{}D(Y,\allowbreak{}A)\allowbreak{}\allowbreak{}=0 for every A\allowbreak{}∈A. They are full subcategories by definition.

Suppose D is triangulated and A is invariant under all shifts. In the triangle \[
 X\xrightarrow fY\xrightarrow gZ\xrightarrow hX[1]
\] the natural exact Hom sequence contains \[
 \operatorname{Hom}(A[1],Z)\longrightarrow\operatorname{Hom}(A,X)
 \xrightarrow{f_*}\operatorname{Hom}(A,Y)
 \longrightarrow\operatorname{Hom}(A,Z).
\] If Z\allowbreak{}∈A\textasciicircum{}⊥,\allowbreak{} both outer groups vanish. Conversely,\allowbreak{} suppose f\_\allowbreak{}* is an isomorphism for every A\allowbreak{}∈A. The exact sequence \[
 \operatorname{Hom}(A,X)\xrightarrow{f_*}\operatorname{Hom}(A,Y)
 \longrightarrow\operatorname{Hom}(A,Z)
 \longrightarrow\operatorname{Hom}(A[-1],X)
 \xrightarrow{f_*}\operatorname{Hom}(A[-1],Y)
\] has its first and last arrows isomorphisms,\allowbreak{} hence Hom(A,\allowbreak{}Z)\allowbreak{}\allowbreak{}=0. This proves 11682–\allowbreak{}11715. The equality notation in its statement means the isomorphism induced by f,\allowbreak{} not an arbitrary abstract group isomorphism.

For the dual statement fix a shift-invariant B. If X\allowbreak{}∈\textasciicircum{}⊥B,\allowbreak{} the exact sequence \[
 \operatorname{Hom}(X[1],B)\longrightarrow\operatorname{Hom}(Z,B)
 \xrightarrow{g^*}\operatorname{Hom}(Y,B)
 \longrightarrow\operatorname{Hom}(X,B)
\] gives the required isomorphism. Conversely,\allowbreak{} apply the assumption to B and B[1]\allowbreak{} in \[
 \operatorname{Hom}(Z,B)\longrightarrow\operatorname{Hom}(Y,B)
 \longrightarrow\operatorname{Hom}(X,B)
 \longrightarrow\operatorname{Hom}(Z[-1],B)
 \longrightarrow\operatorname{Hom}(Y[-1],B).
\] The last arrow is the assumed isomorphism for B[1]\allowbreak{},\allowbreak{} under the shift identifications Hom(Z[−1]\allowbreak{},\allowbreak{}B)\allowbreak{}\allowbreak{}=Hom(Z,\allowbreak{}B[1]\allowbreak{})\allowbreak{}. Thus Hom(X,\allowbreak{}B)\allowbreak{}\allowbreak{}=0. This explicitly proves the omitted dual at 11717–\allowbreak{}11735.

Both orthogonals are strictly full because composition with an isomorphism identifies the relevant Hom groups. They contain zero and are closed under finite biproducts,\allowbreak{} by additivity of Hom. For a retraction Y\allowbreak{}→Z\allowbreak{}→Y,\allowbreak{} the identity composite makes Hom(A,\allowbreak{}Y)\allowbreak{} a retract of Hom(A,\allowbreak{}Z)\allowbreak{},\allowbreak{} and similarly for Hom(Y,\allowbreak{}A)\allowbreak{}. Thus a retract of an orthogonal object remains orthogonal. Shifts preserve orthogonality because Hom(A,\allowbreak{}Y[j]\allowbreak{})\allowbreak{}\allowbreak{}=Hom(A[−j]\allowbreak{},\allowbreak{}Y)\allowbreak{} and the dual formula retain the actual shift isomorphisms. The two exact-sequence criteria show closure under cones; rotation gives two-out-of-three closure. The full triangulated-subcategory criterion at 860–\allowbreak{}880 then supplies the inherited triangulated structure. This proves 11737–\allowbreak{}11758,\allowbreak{} including saturation in the source definition at 1758–\allowbreak{}1763.

The direct sums needed in that proof are finite biproducts. Arbitrary coproduct closure of A\^{}⊥ is not used or inferred. In contrast,\allowbreak{} \^{}⊥A is closed under any existing coproducts by Hom(\allowbreak{}⊕Y\_\allowbreak{}i,\allowbreak{}A)\allowbreak{}\allowbreak{}=∏Hom(Y\_\allowbreak{}i,\allowbreak{}A)\allowbreak{}; A\^{}⊥ is closed under any existing products by Hom(A,\allowbreak{}∏Y\_\allowbreak{}i)\allowbreak{}\allowbreak{}=∏Hom(A,\allowbreak{}Y\_\allowbreak{}i)\allowbreak{}. If every testing object A is compact,\allowbreak{} A\^{}⊥ is also closed under existing coproducts by the defining compact comparison. These precise universal-property statements do not strengthen the source hypotheses silently.

\subsection{2. Existence of approximation triangles is stable under triangles}\label{proofs-11665-12003-existence-of-approximation-triangles-is-stable-under-triangles}

Let A be the source full triangulated subcategory,\allowbreak{} not necessarily strictly full. Suppose X\_\allowbreak{}i has a triangle \[
 A_i\xrightarrow{a_i}X_i\xrightarrow{b_i}B_i
 \xrightarrow{\delta_i}A_i[1],\qquad A_i\in A,\quad B_i\in A^\perp
\] for i\allowbreak{}=1,\allowbreak{}2,\allowbreak{} and let f:\allowbreak{}X\_\allowbreak{}1\allowbreak{}→X\_\allowbreak{}2 be the first map of the given distinguished triangle on X\_\allowbreak{}1,\allowbreak{}X\_\allowbreak{}2,\allowbreak{}X\_\allowbreak{}3. By §1 the map Hom(A\_\allowbreak{}1,\allowbreak{}A\_\allowbreak{}2)\allowbreak{}\allowbreak{}→Hom(A\_\allowbreak{}1,\allowbreak{}X\_\allowbreak{}2)\allowbreak{},\allowbreak{} α\allowbreak{}↦a\_\allowbreak{}2α,\allowbreak{} is an isomorphism. Define α by a\_\allowbreak{}2α\allowbreak{}=fa\_\allowbreak{}1. This specifies the exact commutative square used in the source.

Apply the previously proved nine-object proposition,\allowbreak{} original 1031–\allowbreak{}1052 and PROOFS\_\allowbreak{}0001\_\allowbreak{}1140.md,\allowbreak{} to that square. Choose the middle column to be the specified triangle X\_\allowbreak{}1\allowbreak{}→X\_\allowbreak{}2\allowbreak{}→X\_\allowbreak{}3\allowbreak{}→X\_\allowbreak{}1[1]\allowbreak{},\allowbreak{} and complete the first column on α. The resulting rows are A\_\allowbreak{}i\allowbreak{}→X\_\allowbreak{}i\allowbreak{}→Q\_\allowbreak{}i\allowbreak{}→A\_\allowbreak{}i[1]\allowbreak{},\allowbreak{} and the third column is the distinguished triangle Q\_\allowbreak{}1\allowbreak{}→Q\_\allowbreak{}2\allowbreak{}→Q\_\allowbreak{}3\allowbreak{}→Q\_\allowbreak{}1[1]\allowbreak{}. The lower-right square has the anticommuting sign required by the nine-object proposition; it is not replaced by a commuting square. All other squares commute.

The first two rows are triangles on the original a\_\allowbreak{}i. The triangle morphism axiom gives (id\_\allowbreak{}\{A\_\allowbreak{}i\},\allowbreak{}id\_\allowbreak{}\{X\_\allowbreak{}i\},\allowbreak{}q\_\allowbreak{}i)\allowbreak{} to the originally chosen triangle,\allowbreak{} and the Hom exact sequences show q\_\allowbreak{}i:\allowbreak{}Q\_\allowbreak{}i\allowbreak{}→B\_\allowbreak{}i is an isomorphism. Therefore Q\_\allowbreak{}1,\allowbreak{}Q\_\allowbreak{}2 belong to the strictly full orthogonal,\allowbreak{} and its cone closure gives Q\_\allowbreak{}3\allowbreak{}∈A\textasciicircum{}⊥. The object A\_\allowbreak{}3 is the cone of α. Since A is a full triangulated subcategory,\allowbreak{} there is an object A\_\allowbreak{}3'\allowbreak{}∈A and an isomorphism t:\allowbreak{}A\_\allowbreak{}3'\allowbreak{}→A\_\allowbreak{}3. Transport the third-row triangle along t:\allowbreak{} its first map becomes A\_\allowbreak{}3'\allowbreak{}→\textasciicircum{}tA\_\allowbreak{}3\allowbreak{}→X\_\allowbreak{}3,\allowbreak{} and its last map becomes Q\_\allowbreak{}3\allowbreak{}→A\_\allowbreak{}3[1]\allowbreak{}\allowbreak{}→\textasciicircum{}\{t[1]\allowbreak{}\textasciicircum{}\{-1\}\}A\_\allowbreak{}3'[1]\allowbreak{}. This is precisely a triangle witnessing the source property P(X\_\allowbreak{}3)\allowbreak{},\allowbreak{} without treating A as replete.

The property is invariant under shifts,\allowbreak{} using representatives in A where necessary and the corresponding shift of its triangle with the distinguished-triangle sign convention. Rotating the given triangle therefore proves the other two cases of two-out-of-three. The split triangle X\_\allowbreak{}1\allowbreak{}→X\_\allowbreak{}1\allowbreak{}⊕X\_\allowbreak{}2\allowbreak{}→X\_\allowbreak{}2\allowbreak{}→\textasciicircum{}0X\_\allowbreak{}1[1]\allowbreak{} proves the biproduct case.

For the dual property,\allowbreak{} start with triangles A\_\allowbreak{}i\allowbreak{}→X\_\allowbreak{}i\allowbreak{}→B\_\allowbreak{}i\allowbreak{}→A\_\allowbreak{}i[1]\allowbreak{},\allowbreak{} where B\_\allowbreak{}i\allowbreak{}∈B and A\_\allowbreak{}i\allowbreak{}∈\textasciicircum{}⊥B. The isomorphism Hom(B\_\allowbreak{}1,\allowbreak{}B\_\allowbreak{}2)\allowbreak{}\allowbreak{}→Hom(X\_\allowbreak{}1,\allowbreak{}B\_\allowbreak{}2)\allowbreak{},\allowbreak{} β\allowbreak{}↦βb\_\allowbreak{}1,\allowbreak{} specifies the unique β with βb\_\allowbreak{}1\allowbreak{}=b\_\allowbreak{}2f. Apply the same nine-object construction to the square on X\_\allowbreak{}i\allowbreak{}→B\_\allowbreak{}i; equivalently use the opposite triangulated category,\allowbreak{} whose translation is [−1]\allowbreak{}. Its cone column lies in \^{}⊥B by §1,\allowbreak{} and the B cone is isomorphic to an object of B. Transport along that isomorphism as above. Rotation and the split triangle give all remaining cases. This proves 11813–\allowbreak{}11829 with the actual dual maps.

\subsection{\texorpdfstring{3. Adjoints,\allowbreak{} saturation,\allowbreak{} and double orthogonals}{3.   and double orthogonals}}\label{proofs-11665-12003-adjoints-saturation-and-double-orthogonals}

Let i:\allowbreak{}A\allowbreak{}→D be a full triangulated inclusion. If it has right adjoint v,\allowbreak{} the counit ε\_\allowbreak{}X:\allowbreak{}i vX\allowbreak{}→X induces \[
 \operatorname{Hom}_A(A',vX)
   \xrightarrow{\ \varepsilon_X\circ i(-)\ }
 \operatorname{Hom}_D(iA',X)
\] as the adjunction bijection. Completing ε\_\allowbreak{}X to a triangle and applying §1 puts its cone in A\^{}⊥.

Conversely choose,\allowbreak{} for every X,\allowbreak{} a triangle A\_\allowbreak{}X\allowbreak{}→\textasciicircum{}\{a\_\allowbreak{}X\}X\allowbreak{}→B\_\allowbreak{}X\allowbreak{}→A\_\allowbreak{}X[1]\allowbreak{} as in the source condition (2)\allowbreak{}. Set vX\allowbreak{}=A\_\allowbreak{}X. For f:\allowbreak{}X\allowbreak{}→Y let vf be the unique α with a\_\allowbreak{}Yα\allowbreak{}=fa\_\allowbreak{}X. This uniqueness proves v(id)\allowbreak{}\allowbreak{}=id and v(gf)\allowbreak{}\allowbreak{}=vg vf. The bijections Hom\_\allowbreak{}A(A,\allowbreak{}A\_\allowbreak{}X)\allowbreak{}\allowbreak{}→Hom\_\allowbreak{}D(iA,\allowbreak{}X)\allowbreak{} are induced by a\_\allowbreak{}X,\allowbreak{} and are natural in both A and X by this defining identity. They therefore give the required right adjunction,\allowbreak{} with a\_\allowbreak{}X as counit. The inclusion is fully faithful,\allowbreak{} so the unit A\allowbreak{}→v iA is an isomorphism. Indeed its image under the adjunction is id\_\allowbreak{}\{iA\}; the triangle identities and faithfulness give its two-sided inverse.

If A is strictly full,\allowbreak{} then A⊂\textasciicircum{}⊥(A\textasciicircum{}⊥)\allowbreak{}. For X\allowbreak{}∈\textasciicircum{}⊥(A\textasciicircum{}⊥)\allowbreak{},\allowbreak{} its approximation triangle has b\_\allowbreak{}X\allowbreak{}=0. The split-triangle lemma,\allowbreak{} applied after rotation,\allowbreak{} gives A\_\allowbreak{}X≅X\allowbreak{}⊕B\_\allowbreak{}X[−1]\allowbreak{}. Since B\_\allowbreak{}X[−1]\allowbreak{}\allowbreak{}∈A\textasciicircum{}⊥,\allowbreak{} the projection from A\_\allowbreak{}X to that summand is zero; composing it with the inclusion gives id\_\allowbreak{}\{B\_\allowbreak{}X[−1]\allowbreak{}\}\allowbreak{}=0. Hence B\_\allowbreak{}X\allowbreak{}=0 and a\_\allowbreak{}X is an isomorphism,\allowbreak{} so X\allowbreak{}∈A. Thus A\allowbreak{}=\textasciicircum{}⊥(A\textasciicircum{}⊥)\allowbreak{}. Without strict fullness the same proof gives equality with the strictly full isomorphism closure of A. Saturation in the source\textquotesingle s sense follows:\allowbreak{} any summand X of an object isomorphic to an A object is in \textasciicircum{}⊥(A\textasciicircum{}⊥)\allowbreak{},\allowbreak{} so is isomorphic to an object of A. The other summand is treated by the same argument. No splitting of arbitrary idempotents in D was assumed.

For j:\allowbreak{}B\allowbreak{}→D and a left adjoint w the unit η\_\allowbreak{}X:\allowbreak{}X\allowbreak{}→j wX induces Hom\_\allowbreak{}B(wX,\allowbreak{}B)\allowbreak{}\allowbreak{}→Hom\_\allowbreak{}D(X,\allowbreak{}jB)\allowbreak{},\allowbreak{} β\allowbreak{}↦j(β)\allowbreak{}η\_\allowbreak{}X. Completing η\_\allowbreak{}X backwards gives a triangle A\_\allowbreak{}X\allowbreak{}→X\allowbreak{}→j wX\allowbreak{}→A\_\allowbreak{}X[1]\allowbreak{},\allowbreak{} and §1 places A\_\allowbreak{}X in \^{}⊥B. Conversely,\allowbreak{} such triangles define wX\allowbreak{}=B\_\allowbreak{}X and wf as the unique β with βb\_\allowbreak{}X\allowbreak{}=b\_\allowbreak{}Yf. Uniqueness proves the functor laws and the adjunction naturality. If X\allowbreak{}∈(\textasciicircum{}⊥B)\allowbreak{}\textasciicircum{}⊥,\allowbreak{} then a\_\allowbreak{}X\allowbreak{}=0; the triangle splits as B\_\allowbreak{}X≅X\allowbreak{}⊕A\_\allowbreak{}X[1]\allowbreak{}. The inclusion of this last summand vanishes because Hom(A\_\allowbreak{}X[1]\allowbreak{},\allowbreak{}B\_\allowbreak{}X)\allowbreak{}\allowbreak{}=0,\allowbreak{} so A\_\allowbreak{}X\allowbreak{}=0. Strict fullness gives B\allowbreak{}=(\textasciicircum{}⊥B)\allowbreak{}\textasciicircum{}⊥,\allowbreak{} and without it the equality holds after isomorphism closure. This also proves saturation and the entire omitted dual at 11900–\allowbreak{}11924.

Both adjoints are exact by the full exact-adjoint proof in PROOFS\_\allowbreak{}1844\_\allowbreak{}2401.md §3. Its shift isomorphisms are induced by the actual Hom adjunctions; no assertion of literal equality of chosen objects replaces them.

\subsection{4. The missing hypotheses in the summary proposition}\label{proofs-11665-12003-the-missing-hypotheses-in-the-summary-proposition}

DERIVED-NEW-042 concerns 11954–\allowbreak{}11963. As printed,\allowbreak{} A and B are merely subcategories. Condition (3)\allowbreak{},\allowbreak{} the degree-zero orthogonality and decomposition triangle for every object,\allowbreak{} does not force them to be admissible triangulated subcategories.

An explicit counterexample has D\allowbreak{}=D\textasciicircum{}b(k)\allowbreak{} for a field k,\allowbreak{} A\allowbreak{}=D\textasciicircum{}\{≤0\}(k)\allowbreak{},\allowbreak{} and B\allowbreak{}=D\textasciicircum{}\{≥1\}(k)\allowbreak{},\allowbreak{} both strictly full. Hom\_\allowbreak{}D(A,\allowbreak{}B)\allowbreak{}\allowbreak{}=0:\allowbreak{} use the natural truncation quasi-isomorphisms to a bounded complex supported in degrees ≤0 for A and one supported in degrees ≥1 for B. A bounded complex of vector spaces is a complex of projectives that computes these derived Hom groups; its ordinary chain maps to the second complex are zero by the disjoint degree supports. Transporting through the stated quasi-isomorphisms gives the asserted Hom vanishing for the original A and B.

For the decomposition of an arbitrary bounded complex X,\allowbreak{} use the actual subcomplex A\_\allowbreak{}X with terms X\^{}n for n<0,\allowbreak{} ker(d\_\allowbreak{}X\textasciicircum{}0)\allowbreak{} for n\allowbreak{}=0,\allowbreak{} and zero for n\textgreater0. Its inclusion a\_\allowbreak{}X:\allowbreak{}A\_\allowbreak{}X\allowbreak{}→X has quotient complex B\_\allowbreak{}X\allowbreak{}=X/\allowbreak{}A\_\allowbreak{}X. The quotient has zero terms below zero,\allowbreak{} term X\textasciicircum{}0/\allowbreak{}ker(d\_\allowbreak{}X\textasciicircum{}0)\allowbreak{} in degree zero,\allowbreak{} and terms X\^{}n in degrees n\textgreater0. The degree-zero differential of B\_\allowbreak{}X is injective,\allowbreak{} so H\textasciicircum{}n(B\_\allowbreak{}X)\allowbreak{}\allowbreak{}=0 for n≤0,\allowbreak{} while H\textasciicircum{}n(A\_\allowbreak{}X)\allowbreak{}\allowbreak{}=0 for n\textgreater0. The actual short exact sequence of these complexes gives a distinguished triangle A\_\allowbreak{}X\allowbreak{}→X\allowbreak{}→B\_\allowbreak{}X\allowbreak{}→A\_\allowbreak{}X[1]\allowbreak{}. Thus printed condition (3)\allowbreak{} holds.

But k\allowbreak{}∈A while k[−1]\allowbreak{}\allowbreak{}∉A,\allowbreak{} and k[−1]\allowbreak{}\allowbreak{}∈B while k[−1]\allowbreak{}[1]\allowbreak{}\allowbreak{}=k\allowbreak{}∉B. Neither subcategory is invariant under all shifts,\allowbreak{} so neither is triangulated and (1)\allowbreak{},\allowbreak{}(2)\allowbreak{} both fail. The failure is mathematical,\allowbreak{} not just an omitted word in an otherwise valid proof.

Strict fullness is also not implied by printed (3)\allowbreak{},\allowbreak{} which tests ambient Hom groups and membership of selected objects,\allowbreak{} not the subcategory\textquotesingle s complete morphism sets or isomorphism closure. For example,\allowbreak{} in a nonzero derived category take the subcategory with all objects but only identity and zero maps,\allowbreak{} and let B be its full zero subcategory. The identity triangle X\allowbreak{}→X\allowbreak{}→0 verifies (3)\allowbreak{},\allowbreak{} yet the first subcategory is not full. Hence the conclusion cannot silently restore fullness.

The corrected premise is:\allowbreak{} A and B are \textbf{strictly full subcategories invariant under all shifts}. It is unnecessary to assume their cone closure separately. Here is the complete proof of the missing implication under precisely this premise.

Assume (3)\allowbreak{}. We already have B⊂A\^{}⊥ and A⊂\^{}⊥B. If X\allowbreak{}∈A\textasciicircum{}⊥,\allowbreak{} the first map A\_\allowbreak{}X\allowbreak{}→X of its decomposition triangle vanishes. The split triangle gives B\_\allowbreak{}X≅X\allowbreak{}⊕A\_\allowbreak{}X[1]\allowbreak{}. By shift invariance and (3)\allowbreak{},\allowbreak{} Hom(A\_\allowbreak{}X[1]\allowbreak{},\allowbreak{}B\_\allowbreak{}X)\allowbreak{}\allowbreak{}=0; composing the inclusion of A\_\allowbreak{}X[1]\allowbreak{} with its projection forces its identity to vanish. Thus A\_\allowbreak{}X\allowbreak{}=0 and X≅B\_\allowbreak{}X\allowbreak{}∈B. Since B is strictly full,\allowbreak{} X\allowbreak{}∈B,\allowbreak{} proving B\allowbreak{}=A\textasciicircum{}⊥.

If X\allowbreak{}∈\textasciicircum{}⊥B,\allowbreak{} its map X\allowbreak{}→B\_\allowbreak{}X vanishes,\allowbreak{} hence A\_\allowbreak{}X≅X\allowbreak{}⊕B\_\allowbreak{}X[−1]\allowbreak{}. Shift invariance gives Hom(A\_\allowbreak{}X,\allowbreak{}B\_\allowbreak{}X[−1]\allowbreak{})\allowbreak{}\allowbreak{}=0,\allowbreak{} so the projection to B\_\allowbreak{}X[−1]\allowbreak{} is zero and B\_\allowbreak{}X\allowbreak{}=0. Strict fullness gives X\allowbreak{}∈A. Thus A\allowbreak{}=\textasciicircum{}⊥B. Section 1 now shows that both categories are saturated and triangulated. Section 3 applied to the original decomposition triangles proves both (1)\allowbreak{} and (2)\allowbreak{}. Conversely (1)\allowbreak{} supplies (3)\allowbreak{} by the right-adjoint criterion,\allowbreak{} and (2)\allowbreak{} supplies it by the left-adjoint criterion.

The source patch adds exactly the strict-fullness and shift hypotheses and supplies this implication. It does not assume a missing conclusion in place of proving it.

\subsection{5. Functorial triangles and the exact morphism obstruction}\label{proofs-11665-12003-functorial-triangles-and-the-exact-morphism-obstruction}

DERIVED-UNDER-031 strengthens the statement at 11938–\allowbreak{}11950. For an admissible pair let \[
 A_X\xrightarrow{a_X}X\xrightarrow{b_X}B_X
 \xrightarrow{\delta_X}A_X[1]
\] be the selected triangles,\allowbreak{} with A\_\allowbreak{}X\allowbreak{}∈A,\allowbreak{} B\_\allowbreak{}X\allowbreak{}∈B\allowbreak{}=A\textasciicircum{}⊥. For any f:\allowbreak{}X\allowbreak{}→Y,\allowbreak{} the adjunction gives a unique α\_\allowbreak{}f:\allowbreak{}A\_\allowbreak{}X\allowbreak{}→A\_\allowbreak{}Y with a\_\allowbreak{}Yα\_\allowbreak{}f\allowbreak{}=fa\_\allowbreak{}X. There is also a unique β\_\allowbreak{}f:\allowbreak{}B\_\allowbreak{}X\allowbreak{}→B\_\allowbreak{}Y with β\_\allowbreak{}fb\_\allowbreak{}X\allowbreak{}=b\_\allowbreak{}Yf:\allowbreak{} the relevant exact sequence is \[
 \operatorname{Hom}(A_X[1],B_Y)\longrightarrow
 \operatorname{Hom}(B_X,B_Y)\xrightarrow{b_X^*}
 \operatorname{Hom}(X,B_Y)\longrightarrow\operatorname{Hom}(A_X,B_Y),
\] whose outer terms are zero. The triangle morphism axiom gives a completion of the first square,\allowbreak{} and uniqueness identifies its third map with β\_\allowbreak{}f. Consequently \[
 \delta_Y\beta_f=\alpha_f[1]\delta_X.
\] The unique choices respect identities and composition. In particular,\allowbreak{} two selected triangles on the same X have a \textbf{unique} isomorphism inducing id\_\allowbreak{}X:\allowbreak{} the unique maps in both directions compose to identities by the same uniqueness. These constructions give a natural triangle of functors i v\allowbreak{}→Id\_\allowbreak{}D\allowbreak{}→j w\allowbreak{}→(i v)\allowbreak{}[1]\allowbreak{},\allowbreak{} with the actual exact-adjoint shift identifications. The conclusion is stronger than an unspecified isomorphism between the third objects of two cones.

There is further exact information about arbitrary maps X\allowbreak{}→Y. Using the original a,\allowbreak{}b,\allowbreak{}δ,\allowbreak{} define \[
\begin{aligned}
 d_{-1}(k,l)&=k\delta_X-\delta_Y[-1]\,l,\\
 \phi(h)&=a_Yh b_X,\\
 \rho(f)&=(\alpha_f,\beta_f),\\
 d_0(\alpha,\beta)&=\alpha[1]\delta_X-\delta_Y\beta.
\end{aligned}
\] The domains and codomains form the following sequence,\allowbreak{} exact at each of its three middle terms:\allowbreak{} \[
\begin{gathered}
 \operatorname{Hom}(A_X[1],A_Y)\oplus
 \operatorname{Hom}(B_X,B_Y[-1])
 \xrightarrow{d_{-1}}\operatorname{Hom}(B_X,A_Y)
 \xrightarrow{\phi}\operatorname{Hom}(X,Y)\\
 \xrightarrow{\rho}
 \operatorname{Hom}(A_X,A_Y)\oplus\operatorname{Hom}(B_X,B_Y)
 \xrightarrow{d_0}\operatorname{Hom}(B_X,A_Y[1]).
 \tag{5.1}
\end{gathered}
\] No surjectivity onto the last group is asserted.

For exactness at the pair of Hom groups,\allowbreak{} the displayed triangle morphism identity gives d\_\allowbreak{}0ρ\allowbreak{}=0. Conversely if α[1]\allowbreak{}δ\_\allowbreak{}X\allowbreak{}=δ\_\allowbreak{}Yβ,\allowbreak{} rotate each triangle twice,\allowbreak{} obtaining B\_\allowbreak{}X\allowbreak{}→\textasciicircum{}\{δ\_\allowbreak{}X\}A\_\allowbreak{}X[1]\allowbreak{}\allowbreak{}→\textasciicircum{}\{−a\_\allowbreak{}X[1]\allowbreak{}\}X[1]\allowbreak{}\allowbreak{}→\textasciicircum{}\{−b\_\allowbreak{}X[1]\allowbreak{}\}B\_\allowbreak{}X[1]\allowbreak{} and its Y counterpart. The first square,\allowbreak{} with maps β and α[1]\allowbreak{},\allowbreak{} commutes. The triangle morphism axiom gives a map f[1]\allowbreak{} on the third objects. Its two squares contain the same minus signs on both sides,\allowbreak{} so shifting back gives fa\_\allowbreak{}X\allowbreak{}=a\_\allowbreak{}Yα and b\_\allowbreak{}Yf\allowbreak{}=βb\_\allowbreak{}X. The uniqueness above identifies ρ(f)\allowbreak{}\allowbreak{}=(α,\allowbreak{}β)\allowbreak{}.

If ρ(f)\allowbreak{}\allowbreak{}=0,\allowbreak{} then fa\_\allowbreak{}X\allowbreak{}=0. Exactness of Hom(−,\allowbreak{}Y)\allowbreak{} gives g:\allowbreak{}B\_\allowbreak{}X\allowbreak{}→Y with f\allowbreak{}=g b\_\allowbreak{}X. Since b\_\allowbreak{}Yf\allowbreak{}=0 and Hom(B\_\allowbreak{}X,\allowbreak{}B\_\allowbreak{}Y)\allowbreak{}\allowbreak{}→Hom(X,\allowbreak{}B\_\allowbreak{}Y)\allowbreak{} is injective,\allowbreak{} b\_\allowbreak{}Yg\allowbreak{}=0. The Hom(B\_\allowbreak{}X,\allowbreak{}−)\allowbreak{} sequence for Y gives h:\allowbreak{}B\_\allowbreak{}X\allowbreak{}→A\_\allowbreak{}Y with g\allowbreak{}=a\_\allowbreak{}Yh. Thus f\allowbreak{}=φ(h)\allowbreak{}. Conversely b\_\allowbreak{}Xa\_\allowbreak{}X\allowbreak{}=0 and b\_\allowbreak{}Ya\_\allowbreak{}Y\allowbreak{}=0 show ρφ\allowbreak{}=0 by the defining uniqueness of α and β.

Finally φ(h)\allowbreak{}\allowbreak{}=0 means a\_\allowbreak{}Yh b\_\allowbreak{}X\allowbreak{}=0. The Hom(−,\allowbreak{}Y)\allowbreak{} sequence gives r:\allowbreak{}A\_\allowbreak{}X[1]\allowbreak{}\allowbreak{}→Y with a\_\allowbreak{}Yh\allowbreak{}=rδ\_\allowbreak{}X. Since A\_\allowbreak{}X[1]\allowbreak{}\allowbreak{}∈A,\allowbreak{} its adjunction bijection writes r uniquely as a\_\allowbreak{}Yk for k:\allowbreak{}A\_\allowbreak{}X[1]\allowbreak{}\allowbreak{}→A\_\allowbreak{}Y. Hence a\_\allowbreak{}Y(h−kδ\_\allowbreak{}X)\allowbreak{}\allowbreak{}=0. The inverse-rotated triangle B\_\allowbreak{}Y[−1]\allowbreak{}\allowbreak{}→\textasciicircum{}\{−δ\_\allowbreak{}Y[-1]\allowbreak{}\}A\_\allowbreak{}Y\allowbreak{}→\textasciicircum{}\{a\_\allowbreak{}Y\}Y\allowbreak{}→B\_\allowbreak{}Y gives l:\allowbreak{}B\_\allowbreak{}X\allowbreak{}→B\_\allowbreak{}Y[−1]\allowbreak{} with h−kδ\_\allowbreak{}X\allowbreak{}=−δ\_\allowbreak{}Y[−1]\allowbreak{}l. Therefore h\allowbreak{}=d\_\allowbreak{}\{−1\}(k,\allowbreak{}l)\allowbreak{}. Conversely the triangle composites δ\_\allowbreak{}Xb\_\allowbreak{}X\allowbreak{}=0 and a\_\allowbreak{}Yδ\_\allowbreak{}Y[−1]\allowbreak{}\allowbreak{}=0 give φd\_\allowbreak{}\{−1\}\allowbreak{}=0. This proves every asserted exactness and keeps the rotation sign.

Consequently a compatible pair (α,\allowbreak{}β)\allowbreak{},\allowbreak{} namely d\_\allowbreak{}0(α,\allowbreak{}β)\allowbreak{}\allowbreak{}=0,\allowbreak{} lifts to a map X\allowbreak{}→Y. The set of its lifts is a torsor under the explicitly defined quotient group \[
 \operatorname{Hom}(B_X,A_Y)/\operatorname{Im}(d_{-1}),
\] acting by adding a\_\allowbreak{}Yh b\_\allowbreak{}X. Surjectivity of the action is exactness at Hom(X,\allowbreak{}Y)\allowbreak{},\allowbreak{} and freeness is exactness at Hom(B\_\allowbreak{}X,\allowbreak{}A\_\allowbreak{}Y)\allowbreak{}. Thus the connecting maps δ\_\allowbreak{}X,\allowbreak{}δ\_\allowbreak{}Y determine both the obstruction and the exact ambiguity; one may not replace the category by a direct product merely because each object has two projections. This is an editorial consequence of the original triangles,\allowbreak{} with no novelty claim and no change to their objects or maps.

\subsection{6. The quotient equivalences and their natural isomorphisms}\label{proofs-11665-12003-the-quotient-equivalences-and-their-natural-isomorphisms}

Keep i:\allowbreak{}A\allowbreak{}→D and its exact right adjoint v,\allowbreak{} and j:\allowbreak{}B\allowbreak{}→D and its exact left adjoint w. Then ker v\allowbreak{}=B:\allowbreak{} if vX\allowbreak{}=0 the adjunction makes Hom(A,\allowbreak{}X)\allowbreak{}\allowbreak{}=0 for all A; if X\allowbreak{}∈B then Hom\_\allowbreak{}A(A,\allowbreak{}vX)\allowbreak{}\allowbreak{}=0 for all A,\allowbreak{} so taking A\allowbreak{}=vX shows id\_\allowbreak{}\{vX\}\allowbreak{}=0. Dually ker w\allowbreak{}=A.

Let q\_\allowbreak{}B:\allowbreak{}D\allowbreak{}→D/\allowbreak{}B be the actual Verdier quotient. Since v is exact and vanishes on B,\allowbreak{} its value on every morphism whose cone lies in B is an isomorphism,\allowbreak{} by its triangle exactness. Thus the quotient universal property defines an exact functor \(\bar v:D/B\to A\) with \(\bar v q_B\cong v\). The adjunction unit gives \(\bar v q_B i\cong v i\cong Id_A\). The counit ε\_\allowbreak{}X\allowbreak{}=i vX\allowbreak{}→X has cone B\_\allowbreak{}X\allowbreak{}∈B,\allowbreak{} so q\_\allowbreak{}B(ε\_\allowbreak{}X)\allowbreak{}:\allowbreak{}q\_\allowbreak{}B i vX\allowbreak{}→q\_\allowbreak{}BX is an isomorphism. These isomorphisms are natural for the original arrows by the counit identity; they are natural for inverses of quotient denominators by multiplying that identity by their inverses. They are therefore natural for every roof in D/\allowbreak{}B. This supplies \(q_B i\bar v\cong Id_{D/B}\),\allowbreak{} the other quasi-inverse comparison.

For q\_\allowbreak{}A:\allowbreak{}D\allowbreak{}→D/\allowbreak{}A,\allowbreak{} the exact left adjoint gives \(\bar w:D/A\to B\) with \(\bar w q_A\cong w\). Its counit gives \(\bar w q_A j\cong wj\cong Id_B\). The original unit b\_\allowbreak{}X:\allowbreak{}X\allowbreak{}→j wX has cone A\_\allowbreak{}X[1]\allowbreak{}\allowbreak{}∈A,\allowbreak{} so its quotient isomorphism gives \(Id_{D/A}\cong q_A j\bar w\). Naturality for roofs follows by the same inverse calculation. Thus both quotient equivalences in 11965–\allowbreak{}12003,\allowbreak{} and both descriptions of the adjoints as quotient composites,\allowbreak{} are proved with their actual comparison maps. Object choices give natural isomorphisms,\allowbreak{} not an unproved assertion that independently chosen objects are literally equal.

\subsection{7. Existing reports and corrected receiving quantifiers}\label{proofs-11665-12003-existing-reports-and-corrected-receiving-quantifiers}

P02-E0095 /\allowbreak{} DERIVED-088 report the spelling at original 11744. MC-STK-ERR-0899-OP1 already makes exactly that correction. P02-E0096 /\allowbreak{} DERIVED-089 report the spelling at original 11890; MC-STK-ERR-0900-OP1 already corrects it. Their four physical reports form two reconciliation groups. No additional operation is needed for those spellings.

The source-order receiver search also found the following original reports,\allowbreak{} retained in ADJOINT\_\allowbreak{}RECEIVER\_\allowbreak{}ORIGINAL\_\allowbreak{}REPORTS\_\allowbreak{}PART21.json:\allowbreak{}

\begin{itemize}
\tightlist
\item
  P05-EN-CH36-0056,\allowbreak{} OCC-05101,\allowbreak{} perfect.tex 4848–\allowbreak{}4856.
\item
  P09-ERR-0448,\allowbreak{} OCC-08627,\allowbreak{} spaces-perfect.tex 4459–\allowbreak{}4467.
\item
  P09-ERR-1013,\allowbreak{} OCC-09189,\allowbreak{} spaces-perfect.tex 4459–\allowbreak{}4466.
\end{itemize}

The last two are distinct physical reports of the same source defect. Their original producer/\allowbreak{}ledger identities are preserved,\allowbreak{} and they share one receiving correction group. These reports are separate from the 176-report Derived denominator.

Both receiving passages assert a triangle \[
 E'\longrightarrow E\longrightarrow K\longrightarrow E'[1],
 \qquad E'\in D_{\rm QCoh}(\mathcal O_X),\quad
 \operatorname{Hom}(M,K)=0\quad(M\in D_{\rm QCoh}(\mathcal O_X)).
\] Their universal quantifier incorrectly says ``for every object K''. The arbitrary input in §3 is the middle object,\allowbreak{} so it must say ``for every object E''. K is the cone that is required to lie in the right orthogonal. The subsequent source argument already uses E as its input,\allowbreak{} so changing precisely that quantified letter preserves all following formulas.

The printed statement is demonstrably false as a criterion. Take X\allowbreak{}=Spec(k)\allowbreak{},\allowbreak{} as either a scheme or an algebraic space. The inclusion D\_\allowbreak{}QCoh(O\_\allowbreak{}X)\allowbreak{}\allowbreak{}→D(O\_\allowbreak{}X)\allowbreak{} is the identity and has a right adjoint. But demanding the displayed orthogonality for every K includes K\allowbreak{}=k and M\allowbreak{}=k,\allowbreak{} for which Hom(k,\allowbreak{}k)\allowbreak{} contains id\_\allowbreak{}k≠0. This validates the producer reports mathematically,\allowbreak{} rather than matching their loci only.

For the corrected scheme argument,\allowbreak{} the source\textquotesingle s Mayer–\allowbreak{}Vietoris triangle has middle-input objects Rj\_\allowbreak{}\{U,\allowbreak{}*\}E|\_\allowbreak{}U,\allowbreak{} Rj\_\allowbreak{}\{V,\allowbreak{}*\}E|\_\allowbreak{}V,\allowbreak{} Rj\_\allowbreak{}\{U∩V,\allowbreak{}*\}E|\_\allowbreak{}\{U∩V\}. The preparation lemma in §2 propagates approximation existence through that triangle and its finite sum. For one of these open maps,\allowbreak{} start with the triangle L'\allowbreak{}→L\allowbreak{}→Q\allowbreak{}→L'[1]\allowbreak{} on U. Applying the actual exact Rj\_\allowbreak{}* gives Rj\_\allowbreak{}*L'\allowbreak{}→Rj\_\allowbreak{}*L\allowbreak{}→Rj\_\allowbreak{}*Q\allowbreak{}→Rj\_\allowbreak{}*L'[1]\allowbreak{}. The cited quasi-coherent pushforward result puts Rj\_\allowbreak{}*L' in the required subcategory. Its cone lies in the orthogonal because \[
 \operatorname{Hom}_X(M,Rj_*Q)
   \cong\operatorname{Hom}_U(Lj^*M,Q)=0,
\] with the actual pullback preserving the quasi-coherent category. The algebraic-space passage uses the same maps for the stated étale morphisms and the elementary distinguished-square triangle; its pullback–\allowbreak{}pushforward adjunction is the same precise Hom comparison. The quantifier correction therefore repairs both receiving uses. Their separate geometric dependencies remain in their own chapter reviews; this bounded check does not certify those entire chapters.

\subsection{8. Other receiving uses and limits of this review}\label{proofs-11665-12003-other-receiving-uses-and-limits-of-this-review}

In stacks-perfect.tex 861–\allowbreak{}934 the kernel category \[
 B=D_{\textit{Parasitic}\cap\textit{LQCoh}^{fbc}}(\mathcal O_X)
 \subset D_{\textit{LQCoh}^{fbc}}(\mathcal O_X)
\] is the stated full derived subcategory and is invariant under all shifts. Its left orthogonal A is strictly full and shift invariant by §1. Thus the pair in the source remark satisfies the repaired premise of DERIVED-NEW-042. The preceding lemma provides E'\allowbreak{}=Lg\_\allowbreak{}!g\textasciicircum{}*E\allowbreak{}→E with cone P\allowbreak{}∈B and the actual comparison \[
 \operatorname{Hom}(Lg_!g^*E,P')
  \cong\operatorname{Hom}(g^*E,g^*P')=0.
\] It supplies the original decomposition and orthogonality,\allowbreak{} so the admissibility and quotient conclusions follow from §§3,\allowbreak{}4,\allowbreak{}6 with no new change to this receiver. Section 5 further proves that its canonical triangle is unique up to a unique isomorphism fixing E and is functorial in E. That strengthening is recorded editorially. The complete independent review of the preceding stack/\allowbreak{}topos comparison lemmas remains in their own chapter queues.

In equiv.tex 2821–\allowbreak{}2864 the essential image A of the fully faithful exact F is taken with all its ambient morphisms and all objects isomorphic to its image; it is strictly full. Exactness preserves shifts. Full faithfulness lifts every map between image objects to the source; applying F to its source cone proves cone closure up to the prescribed isomorphism. Thus A is triangulated. If F\_\allowbreak{}r is its right adjoint as cited there,\allowbreak{} the counit F F\_\allowbreak{}rN\allowbreak{}→N gives the right approximation:\allowbreak{} for a source object M,\allowbreak{} adjunction and the fully faithful unit identify Hom(FM,\allowbreak{}F F\_\allowbreak{}rN)\allowbreak{}\allowbreak{}→Hom(FM,\allowbreak{}N)\allowbreak{} with an isomorphism. Section 3 gives the right admissibility and A\allowbreak{}=\textasciicircum{}⊥(A\textasciicircum{}⊥)\allowbreak{} used by the receiver. Its later point-sheaf detection argument is not a new condition needed for this categorical step. It remains subject to its own source review.

The root-chapter reference search was bounded to explicit references to this section\textquotesingle s labels; it is not a claim of exhaustive semantic dependency discovery. Later Postnikov source was read ahead through 12150 only,\allowbreak{} and is not covered by this completed review.
